\documentclass[10pt,a4paper]{amsart}
\usepackage[margin=19mm]{geometry}
\usepackage{amsmath,amssymb,mathtools}
\usepackage[scr=rsfs,cal=euler]{mathalfa}
\usepackage{xfrac}
\usepackage[unicode,hidelinks,bookmarks=false]{hyperref}
\newcommand{\ie}{i.e.\ }
\newcommand{\cf}{cf.\ }
\newcommand{\resp}{respectively\ }
\newcommand{\Subsection}{\subsection}
\providecommand{\nobreakdash}{\nobreak\hbox{-}}
\setlength{\emergencystretch}{2em}
\multlinegap0pt
\usepackage{bm}
\usepackage{bbm}
\usepackage{dsfont}
\usepackage{xspace}
\usepackage{cancel}
\usepackage{mathtools}
\usepackage{amsthm}
\makeatletter
\@ifundefined{lemma}{\newtheorem{lemma}{Lemma}}{}
\makeatother
\title[Trees with maximum $\sigma$-irregularity under a prescribed ]{Trees with maximum $\sigma$-irregularity under a prescribed maximum degree 6}
\author{Milan Ba\v{s}i\'c}
\date{Source: arXiv:2602.01262v1 (CC BY 4.0)}
\begin{document}
\maketitle

\noindent\emph{Source and licence.} Trees with maximum $\sigma$-irregularity under a prescribed maximum degree 6. Version arXiv:2602.01262v1 (CC BY 4.0), distributed under the Creative Commons Attribution 4.0 International licence, as recorded in the arXiv licence metadata for this exact version and shown on the abstract page https://arxiv.org/abs/2602.01262v1. Full rights notice: licenses/arxiv-2602.01262-CC-BY-4.0.txt.\par\smallskip

\noindent\textit{Editorial scope.} Counted unit: the Lemma whose source label is \texttt{lem:no4-threshold} in the article (source0.tex lines 548--557, source label \texttt{lem:no4-threshold}), delivered together with the article's own complete original proof of it (source0.tex lines 559--741, one \texttt{proof} environment of 183 lines). No mathematics is written, repaired or re-derived: statement and proof are the article's own text line for line. Required context retained uncounted: eq:tree-ni-6 (lines 289--298); tab:F-delta6-row (lines 334--353); lem:no5-threshold (lines 457--464). Every definition, display and label of those lines is retained unchanged. Omitted from this package and not redistributed: 1--288, 299--333, 354--456, 465--547, 558--558, 742--2931. Nothing inside a retained excerpt is omitted or altered: every retained body is delivered exactly as the article's lines are written, so no omission receipt of any kind accompanies this package. Citations: the retained excerpts carry no citation command at all, so no bibliography entry of the article is transcribed and no reference is redistributed. Reference adaptation: none was needed and none was made; every reference inside the delivered unit resolves inside it, so no unresolved reference or citation is produced. Printed slips kept literal: no printed word, symbol, hypothesis or number of the counted statement or of its proof was altered. Layout and class: the article's own class is \texttt{elsarticle} with packages including natbib, graphicx, multirow, amsmath, amssymb, amsfonts, amsthm, mathrsfs, appendix, xcolor, textcomp, manyfoot, booktabs, algorithm; the wrapper is the lane's pinned \texttt{amsart} editorial wrapper at 10pt on A4, which restates the theorem-like environments the retained text needs and adds the article's own macro definitions that the retained text needs (none), with extra wrapper packages (bm, bbm, dsfont, xspace, cancel, mathtools). The delivered numbering is the unit's own. Assets: no retained span contains a figure environment, a graphics-inclusion command, a PDF-page inclusion, a table or a TikZ picture; no image file is copied into this package, the package declares no asset dependency and no image file is redistributed with it. Licence: version arXiv:2602.01262v1 of this article is distributed under the Creative Commons Attribution 4.0 International licence, as recorded in the arXiv licence metadata for this exact version and shown on the abstract page https://arxiv.org/abs/2602.01262v1. A full rights notice is delivered at licenses/arxiv-2602.01262-CC-BY-4.0.txt. Characters: every retained excerpt is pure ASCII, so the delivered engine, XeLaTeX with its default Latin Modern text font, reports no missing character.
\begin{align}
n_1+n_2+\cdots+n_6&=n, \label{eq:tree-sum-ni}\\
n_1+2n_2+\cdots+6n_6&=2n-2, \label{eq:tree-sum-deg}\\
m_{12}+m_{13}+m_{14}+m_{15}+m_{16}&=n_1, \label{eq:tree-ni-1}\\
m_{12}+2m_{22}+m_{23}+m_{24}+m_{25}+m_{26}&=2n_2, \label{eq:tree-ni-2}\\
m_{13}+m_{23}+2m_{33}+m_{34}+m_{35}+m_{36}&=3n_3, \label{eq:tree-ni-3}\\
m_{14}+m_{24}+m_{34}+2m_{44}+m_{45}+m_{46}&=4n_4, \label{eq:tree-ni-4}\\
m_{15}+m_{25}+m_{35}+m_{45}+2m_{55}+m_{56}&=5n_5, \label{eq:tree-ni-5}\\
m_{16}+m_{26}+m_{36}+m_{46}+m_{56}+2m_{66}&=6n_6. \label{eq:tree-ni-6}
\end{align}

\begin{table}[ht]
\centering
\caption{Values of the penalty function $F(i,j)$ for $\Delta=6$, arranged row-wise
in increasing order for $(i,j)\in S=\{(i,j):1\le i\le j\le 6,\ (i,j)\notin\{(1,6),(2,6)\}\}$.}
\label{tab:F-delta6-row}
\begin{tabular}{cccc}
\hline
$(i,j)$ & $F(i,j)$ & $(i,j)$ & $F(i,j)$ \\ \hline
$(3,6)$ & $4.0$   & $(4,6)$ & $7.5$   \\
$(2,5)$ & $7.6$   & $(1,5)$ & $9.6$   \\
$(3,5)$ & $9.6$   & $(5,6)$ & $9.6$   \\
$(6,6)$ & $10.0$  & $(4,5)$ & $11.1$  \\
$(5,5)$ & $11.2$  & $(4,4)$ & $13.0$  \\
$(3,4)$ & $13.5$  & $(2,4)$ & $13.5$  \\
$(3,3)$ & $16.0$  & $(1,4)$ & $17.5$  \\
$(2,3)$ & $18.0$  & $(2,2)$ & $22.0$  \\
$(1,3)$ & $24.0$  & $(1,2)$ & $30.0$  \\
$(1,1)$ & $40.0$  &         &         \\ \hline
\end{tabular}
\end{table}

\begin{lemma}
\label{lem:no5-threshold}
Let $\Delta=6$ and let $T$ be a tree.
\begin{itemize}
\item If $P(T)\le 30$, then $n_5=0$.
\item If $n\not\equiv 0\pmod 6$ and $P(T)\le 40$, then $n_5=0$.
\end{itemize}
\end{lemma}

\begin{lemma}
\label{lem:no4-threshold}
Let $\Delta=6$ and let $T$ be a tree on $n$ vertices.

\begin{itemize}
\item If $P(T)<30$, then $n_4=0$.
\item If $n\not\equiv 1 \pmod 6$ and $P(T)=30$, then $n_4=0$.
\item If $n\not\equiv \{0,1,2\} \pmod 6$ and $30<P(T)\le 40$, then $n_4=0$.
\end{itemize}
\end{lemma}

\begin{proof}
Assume $n_4\ge 1$.

From the degree--$4$ handshake equation we have
\[
m_{1,4}+m_{2,4}+m_{3,4}+m_{4,5}+m_{4,6}+2m_{4,4}=4n_4 \ge 4 .
\]
Hence there are at least $4$ incidences at degree--$4$ vertices. 
Indeed, if $m_{4,4}=0$, this is immediate from $4n_4\ge 4$, while if
$m_{4,4}\ge 1$, then necessarily $n_4\ge 2$ and hence $4n_4\ge 8$.
By Table~\ref{tab:F-delta6-row}, the smallest penalty among all edge--types
incident to degree $4$ is $F(4,6)=7.5$, so
\[
P(T)\ \ge\ 4\cdot 7.5 \ =\ 30 .
\]
This proves the first item.

\medskip
\noindent
\textbf{Case 1: $P(T)=30$ and $n_4\ge 1$.}  
Equality in the above bound forces all incidences at degree--$4$ vertices
to come from $(4,6)$--edges. Hence
\[
n_4=1, \qquad m_{4,6}=4, \qquad m_{i,j}=0 \ \text{for all }(i,j)\in S\setminus\{(4,6)\}.
\]
Let $x=m_{1,6}$ and $y=m_{2,6}$.  
Counting edges gives
\[
x+y+4=n-1 ,
\]
while the degree--$6$ handshake equation (\ref{eq:tree-ni-6}) yields
\[
x+y+4=6n_6 .
\]
Thus $6n_6=n-1$, i.e.\ $n\equiv 1\pmod 6$.  
Therefore, if $n\not\equiv 1\pmod 6$, then $P(T)=30$ is impossible with $n_4\ge 1$,
which proves the second item.

\medskip
\noindent
\textbf{Case 2: $30<P(T)\le 40$ and $n\not\equiv\{0,1,2\}\pmod 6$.}  
Assume first that $n_4\ge 2$. Then
\[
m_{1,4}+m_{2,4}+m_{3,4}+m_{4,5}+m_{4,6}+2m_{4,4}=4n_4 \ge 8 .
\]

If $m_{4,4}=0$, then all incidences come from edges incident to degree $4$
of types $(1,4)$, $(2,4)$, $(3,4)$, $(4,5)$ or $(4,6)$.  
The smallest penalty among these is $F(4,6)=7.5$, hence
\[
P(T)\ \ge\ 8\cdot 7.5 \ =\ 60 \ >\ 40 ,
\]
a contradiction.

If $m_{4,4}\ge 1$, then one $(4,4)$--edge contributes two incidences at cost $13$,
and the remaining at least six incidences contribute at least $6\cdot 7.5=45$.
Hence
\[
P(T)\ \ge\ 13 + 45 \ =\ 58 \ >\ 40 ,
\]
again a contradiction.

Therefore $n_4\ge 2$ is impossible, and we must have $n_4=1$.

By Lemma~\ref{lem:no5-threshold} (applied under $P(T)\le 40$ and
$n\not\equiv 0\pmod 6$), we also have $n_5=0$, hence $m_{4,5}=0$ and $m_{4,4}=0$.
Thus the degree--$4$ handshake reduces to
\[
m_{1,4}+m_{2,4}+m_{3,4}+m_{4,6}=4 .
\]


\medskip
\noindent\textbf{Case 2.1: $m_{4,6}\le 2$.}
Then $m_{1,4}+m_{2,4}+m_{3,4}\ge 2$, so at least two incidences come from
$(1,4)$, $(2,4)$, or $(3,4)$.
By Table~\ref{tab:F-delta6-row}, each such incidence has penalty at least $13.5$.
Hence
\[
P(T)\ge 2\cdot 7.5+2\cdot 13.5=42>40,
\]
a contradiction. Therefore $m_{4,6}\in\{3,4\}$.


\medskip
\noindent
\textbf{Case 2.2: $m_{4,6}=4$.}  
Then $m_{1,4}=m_{2,4}=m_{3,4}=0$, and the remaining penalty budget is at most
$P(T)-30\le 10$.  
Since $F(3,6)=4$ and $F(6,6)=10$ are the only values in $S$ not exceeding $10$
(with $n_5=0$), we must have
\[
m_{i,j}=0 \ \text{for all }(i,j)\in S\setminus\{(3,6),(4,6),(6,6)\}, \qquad
m_{3,6}\le 2, \quad m_{6,6}\le 1 .
\]

Let $x=m_{1,6}$, $y=m_{2,6}$, $t=m_{3,6}$ and $b=m_{6,6}$.  
Counting edges gives
\begin{eqnarray}
\label{edge_count_4}
  x+y+t+b+4=n-1 ,  
\end{eqnarray}

and the degree--$6$ handshake equation yields
\begin{eqnarray}
\label{handshake_6_4}
x+y+t+4+2b=6n_6 .
\end{eqnarray}
Subtracting (\ref{edge_count_4}) from (\ref{handshake_6_4}) gives
\[
b=6n_6-(n-1) .
\]
Since $n\not\equiv\{0,1\}\pmod 6$, we obtain $b\ge 2$, contradicting $b\le 1$.
Hence this case is impossible.

\medskip
\noindent
\textbf{Case 2.3: $m_{4,6}=3$.}  
Then
\[
m_{1,4}+m_{2,4}+m_{3,4}=1 .
\]

\smallskip
\noindent
\textbf{Case 2.3(a): the unique non-$(4,6)$ edge is $(1,4)$.}  
Then
\[
P(T)\ \ge\ 3\cdot 7.5 + 17.5 \ =\ 40 ,
\]
so $P(T)=40$ and no further $(i,j)\in S$ may occur.  
Hence $m_{6,6}=m_{3,6}=0$.

Let $x=m_{1,6}$ and $y=m_{2,6}$.  
Edge counting gives
\[
x+y+3+1=n-1 \quad\Rightarrow\quad x+y=n-5 ,
\]
while the degree--$6$ handshake gives
\[
x+y+3=6n_6 .
\]
Thus $6n_6=n-2$, impossible for $n\not\equiv 2\pmod 6$.

\smallskip
\noindent
\textbf{Case 2.3(b): the unique non-$(4,6)$ edge is $(2,4)$ or $(3,4)$.}  
Then
\[
P(T)\ \ge\ 3\cdot 7.5 + 13.5 \ =\ 36 ,
\]
so the remaining budget is at most $4$.  
Since $F(3,6)=4$ is the only admissible positive value in $S$ not exceeding $4$,
we must have
\[
m_{i,j}=0 \ \text{for all }(i,j)\in S\setminus\{(3,6),(4,6),(2,4),(3,4)\},
\qquad m_{3,6}\le 1 .
\]

The degree--$3$ handshake equation becomes
\[
m_{3,4}+m_{3,6}=3n_3 .
\]
Since $m_{3,4}\in\{0,1\}$ and $m_{3,6}\le 1$, this forces
\[
m_{3,4}=m_{3,6}=0, \qquad n_3=0 .
\]

Now edge counting gives
\[
x+y+3+1=n-1 \quad\Rightarrow\quad x+y=n-5 ,
\]
and the degree--$6$ handshake yields
\[
x+y+3=6n_6 .
\]
Thus again $6n_6=n-2$, impossible for $n\not\equiv 2\pmod 6$.

\medskip
Since all cases lead to contradictions, we conclude that $n_4=0$ whenever
$n\not\equiv\{0,1,2\}\pmod 6$ and $30<P(T)\le 40$.  
This proves the third item.
\end{proof}

\end{document}
